\documentclass[11pt]{article}

%%METADATA
\title{GA1025 Macro Theory \\Solution to Problem Set 4}
\author{
Junbiao Chen\thanks{E-mail: jc14076@nyu.edu.}
}
\date{\today}


%%PACKAGES
\usepackage{mdframed} % For boxed environments
\usepackage{graphicx}
\usepackage{grffile}
\usepackage{tabularx}
\usepackage{setspace}
\usepackage{amsmath,amsthm,amssymb}
\usepackage[hyphens]{url}
\usepackage{natbib}
\usepackage[font=normalsize,labelfont=bf]{caption}
\usepackage[margin=1in]{geometry}
\usepackage{hyperref}
\hypersetup{colorlinks=true,urlcolor=blue,citecolor=blue}
\usepackage{stmaryrd}  %Package with \boxast command
\usepackage{enumerate}% http://ctan.org/pkg/enumerate %Supports lowercase Roman-letter enumeration
\usepackage{verbatim} %Package with \begin{comment} environment
%\usepackage{enumitem}
\usepackage{physics}
\usepackage{tikz}
\usepackage{listings}
\usepackage{upquote}
\usepackage{booktabs} %Package with \toprule and \bottomrule
\usepackage{etoc}     %Package with \localtableofcontents
\usepackage{placeins}    %Package that prevent repositioning the tables
\usepackage{multicol}
\usepackage{bm}
\usepackage{subfig}
\usepackage{csquotes}

\definecolor{dkgreen}{rgb}{0,0.6,0}
\definecolor{gray}{rgb}{0.5,0.5,0.5}
\definecolor{mauve}{rgb}{0.58,0,0.82}

\lstset{language=bash,
  frame=tb,
  aboveskip=3mm,
  belowskip=3mm,
  showstringspaces=false,
  columns=flexible,
  basicstyle={\small\ttfamily},
  numbers=none,
  numberstyle=\tiny\color{gray},
  keywordstyle=\color{blue},
  commentstyle=\color{dkgreen},
  stringstyle=\color{mauve},
  breaklines=true,
  breakatwhitespace=false,
  tabsize=3
}

\lstset{language=C,
  aboveskip=3mm,
  belowskip=3mm,
  showstringspaces=false,
  columns=flexible,
  basicstyle={\small\ttfamily},
  numbers=none,
  numberstyle=\tiny\color{gray},
  keywordstyle=\color{blue},
  commentstyle=\color{dkgreen},
  stringstyle=\color{mauve},
  breaklines=true,
  breakatwhitespace=false,
  tabsize=4
}

\definecolor{lightblue}{rgb}{0.68, 0.85, 0.9} %

%CUSTOM DEFINITIONS
\theoremstyle{definition}
\newtheorem{definition}{Definition}[section]
\newtheorem*{remark}{Remark}
\setcounter{secnumdepth}{3}
\usepackage{tikz}
\usetikzlibrary{arrows.meta}
\usetikzlibrary{automata, positioning, arrows, calc}

\tikzset{
	->,  % makes the edges directed
	>=stealth, % makes the arrow heads bold
	shorten >=2pt, shorten <=2pt, % shorten the arrow
	node distance=3cm, % specifies the minimum distance between two nodes. Change if n
	every state/.style={draw=blue!55,very thick,fill=blue!20}, % sets the properties for each ’state’ n
	initial text=$ $, % sets the text that appears on the start arrow
 }

%% PROPOSITION
% Define the Proposition environment
\newmdenv[
  innerleftmargin=10pt, 
  innerrightmargin=10pt,
  innertopmargin=10pt,
  innerbottommargin=10pt,
  linecolor=black, 
  linewidth=1pt,
  backgroundcolor=white, 
  roundcorner=5pt
]{propositionbox}

\newtheoremstyle{boldtitle} % Define a new theorem style
  {10pt} % Space above
  {10pt} % Space below
  {\itshape} % Body font
  {} % Indent amount
  {\bfseries} % Theorem head font
  {.} % Punctuation after theorem head
  { } % Space after theorem head
  {} % Theorem head spec

\theoremstyle{boldtitle} % Use the custom style
\newtheorem{proposition}{Proposition} % Define the proposition environment

% Redefine the proposition environment to use the box
\newenvironment{boxedproposition}[1][]
{\begin{propositionbox}\begin{proposition}[#1]}
{\end{proposition}\end{propositionbox}}

%%FORMATTING
\usepackage[bottom]{footmisc}
\onehalfspacing
\numberwithin{equation}{section}
\numberwithin{figure}{section}
\numberwithin{table}{section}
\bibliographystyle{../bib/aeanobold-oxford}


%main text
\begin{document}

\maketitle

\section*{Exercise 2.1 from RMT5}
\paragraph{Solution} 
For the Markov chain 
\[
(\mathbf{P}, \pi_0) = 
\left(\left[
\begin{matrix}
  0.9 &  0.1 \\ 0.3 &  0.7  
\end{matrix}
\right], 
\left[
\begin{matrix}
    0.5 \\ 0.5
\end{matrix}
\right]
\right)
\]
random variable $y_t = \bar{y}' x_t$, where $\bar{y}' = [1,5]$.

The likelihood function of the data is given by 
\begin{itemize}
    \item(a) $L = 0.5 P_{12}^2 P_{21}^2 = 0.5 \cdot 0.1^2 \codt 0.3^2$
    \item(b) $L = 0.5 P_{11}^4 = 0.5 \cdot 0.9^4$
    \item(c) $L = 0.5 P_{22}^4 = 0.5 \cdot 0.7^4$
\end{itemize}

\section*{Exercise 2.2 from RMT5}
\subsection*{(a)} Assume the transition probability matrix is 
\[
\left[
\begin{matrix}
  p&  (1-p) \\ (1-q) &  q 
\end{matrix}
\right]
\]
Then we have 
\begin{align*}
    p + 5(1-p) & = 1.8 \\
    (1 - q) + 5q & = 3.4
\end{align*}
Solving them, we have $p = 0.8$ and $q = 0.6$.

\subsection*{(a)} In this problem, we use the cross-sectional data to compute the prabability matrix.
Assume the transition probability matrix is 
\[
\left[
\begin{matrix}
  p&  (1-p) \\ (1-q) &  q 
\end{matrix}
\right]
\]
Then we have 
\begin{align*}
    p + 5(1-p) & = 1.8 \\
    p^2 + 5 p(1-p) + (1-p)(1-q) + 5(1-p)q & = 2.12 \\
\end{align*}
Solving them, we have $p = 0.8$ and $q = 0.6$.

\section*{Exercise 2.3 from RMT5}
\paragraph{Q1 Solution}
Note that $\mathbb{E}(y_t|x_0) = \mathbf{P}^t \bar{y}$.
Denote $u(c_t)$ by $u_t$.
Then we have 
\begin{align*}
    \bar{v} & = \mathbb{E}(\sum_{t=0}^{\infty}\beta^t u_t |x_0) \\ 
        & = \sum_{t=0}^{\infty}\beta^t  \mathbb{E}(u_t |x_0) \\ 
        & = \sum_{t=0}^{\infty}\beta^t  \mathbf{P}^t \bar{u} \\ 
        & = (I - \beta  \mathbf{P})^{-1} \bar{u}
\end{align*}
By the law of total probability, we can express the unconditional expected value starting at time 0 as:
\begin{align*}
    V & = \mathbb{E}[v_0] \\ 
     & = \mathbb{E}[\bar{v}' x_0] \quad \text{ (by definition)}\\
     & = \sum_{i=1}^{n} \bar{v}_i \text{Pr}(x_0 = e_i) \\
     & = \sum_{i=1}^{n} \bar{v}_i \pi_{0,i} 
\end{align*}

\paragraph{Q2 Solution}
Based on the results in Q1, we only need to compute $\pi_0' (I - \beta P)^{-1}$ and then compare
Since 
\[
\left[
\begin{matrix}
  1-\beta & 0 \\ 0 &  1 - \beta
\end{matrix}
\right]^{-1} 
= 
\left[
\begin{matrix}
  1-\frac{1}{2}\beta & -\frac{1}{2}\beta \\ -\frac{1}{2}\beta &  1 -\frac{1}{2}\beta
\end{matrix}
\right]^{-1} 
\]
We can conclude that processes 1 and 2 yeild the same unconditional expected value for both 
$\gamma = 2.5$ and $4$.


%% =========== %% ========= %%
\bibliography{../bib/references.bib}

\end{document}